% ch9_f 第9章 §9.5–§9.6 (书页 419–430 / p434–p445)
% 块边界判定：
%   起点——书页 419（p434）首段 "Knowing the translational symmetry ..." 为完整新句；
%     前块止于书页 418 末 "The gapless spinons appear only at (pi/2, +/-pi/2)."，
%     无前块溢出内容，故本块不设 %JOIN%，亦无需 TAIL 区。
%   实际覆盖——书页 419 起于 9.6.1 节后部（式 (9.6.9) 之前的段落），依次为 9.6.2 节、
%     9.7 节、9.8 节（含 9.8.1、9.8.2）、9.9 节引言；9.5 节与 9.6/9.6.1 节标题、
%     式 (9.6.1)--(9.6.8)、图 9.8、图 9.9 均在前块（书页 415--418）内。
%   终点——书页 430（p445）末句 "... in the light of quantum orders and their PSG
%     characterization." 为 9.9 节引言末句，书页 431（p446）顶部起新小节 9.9.1，
%     本块自然收尾，不设 %--E-TAIL--。
% 蓝本问题记录：
%   1) 问题 9.6.3(a) 原书把 eta/chi 与两个态的对应关系印反（与 9.6.2 节正文矛盾：
%      正文为 chi=0 -> SU2An0、eta=0 -> SU2Bn0，与拟设形式逐项比对一致），
%      且将 SU2Bn0 误标为 "SU(2)-gapless"，已按正文与本意订正并加译注。
%   2) 式 (9.8.9) 末行原书漏印下标 0（作 $a^{1,2}=0$），按本意作 $a_0^{1,2}=0$。
%   3) 式 (9.6.9) 原书即印作 $k_1+k_2+\pi x$（而变换式为 $|\pmb{k}\rangle\to
%      |\pmb{k}+\pi\pmb{y}\rangle$），照原图转写。
%   4) 图 9.10 图注 "the two-spinon spectrum docs have" 中 "docs" 为 "does" 之误。

%FIG{9.10}: {$Z_2$ 自旋液体自旋子色散 $E_+(\pmb{k})$ 作为 $(k_x/2\pi,\,k_y/2\pi)$ 的函数的等高线图：(a) 式 (9.2.39) 中的 $Z_2$ 无能隙态 Z2Ax2(12)n；(b) 式 (9.6.10) 中的 $Z_2$ 二次方态 Z2Bx2(12)n。尽管 (a) 中的单自旋子色散缺乏旋转对称性和宇称对称性，双自旋子谱却具有这些对称性。}

虽然上述 Z2B 自旋液体具有平移对称性，但发现自旋子谱只在格点布里渊区的一半上有定义，这一点看起来很奇怪。不过，这与物理自旋液体中的平移对称性并不矛盾，因为单自旋子激发并不是物理的。只有双自旋子激发才对应于物理激发，其谱应当定义在满布里渊区上。现在的问题是：如何从只在半布里渊区上有定义的单自旋子谱，得到定义在满布里渊区上的双自旋子谱。设 $|\pmb{k},1\rangle$ 与 $|\pmb{k},2\rangle$ 是单自旋子的两个本征态，能量为正，分别为 $E_1(\pmb{k})$ 与 $E_2(\pmb{k})$（这里 $k_x\in(-\pi/2,\pi/2)$，$k_y\in(-\pi,\pi)$）。沿 $\pmb{x}$ 的平移（随后再作一个规范变换）把 $|\pmb{k},1\rangle$ 与 $|\pmb{k},2\rangle$ 变成能量相同的另外两个本征态：
\begin{equation*}
\begin{aligned}
|\pmb{k},1\rangle &\to |\pmb{k}+\pi\pmb{y},1\rangle \\
|\pmb{k},2\rangle &\to |\pmb{k}+\pi\pmb{y},2\rangle
\end{aligned}
\end{equation*}
于是我们看到，双自旋子态 $|\pmb{k}_1,\alpha_1\rangle|\pmb{k}_2,\alpha_2\rangle\pm|\pmb{k}_1+\pi\pmb{y},\alpha_1\rangle|\pmb{k}_2+\pi\pmb{y},\alpha_2\rangle$ 的动量与能量由下式给出：
\begin{equation*}
\begin{aligned}
E_{2\text{-spinon}} &= E_{\alpha_1}(\pmb{k}_1)+E_{\alpha_2}(\pmb{k}_2) \\
\pmb{k} &= \pmb{k}_1+\pmb{k}_2,\ \ \pmb{k}_1+\pmb{k}_2+\pi\pmb{x}
\end{aligned}
\tag{9.6.9}
\end{equation*}
式 (9.6.9) 使我们能够由单自旋子谱构造出双自旋子谱。

\subsection{$SU(2)$ 自旋液体附近的一个奇怪的对称自旋液体}

在式 (9.2.32) 的均匀 RVB 态（即 $SU(2)$ 无能隙态 SU2An0）附近以及式 (9.2.33) 的 $\pi$ 通量态（即 $SU(2)$ 线性态 SU2Bn0）附近，存在许多类型的对称拟设。这里我们只考虑其中的一个——$Z_2$ 自旋液体 Z2By1(12)n（注意 Z2By1(12)n 与 Z2Bx2(12)n 规范等价），它具有如下形式：
\begin{equation*}
u_{\pmb{i},\pmb{i}+\pmb{x}}=\mathrm{i}\chi\tau^{0}+\eta\tau^{1},\qquad
u_{\pmb{i},\pmb{i}+\pmb{y}}=(-)^{i_x}(\mathrm{i}\chi\tau^{0}+\eta\tau^{2}),\qquad
a_0^{1,2,3}=0.
\tag{9.6.10}
\end{equation*}
上述拟设在 $\chi=0$ 时退化为 $SU(2)$ 无能隙态 SU2An0 的拟设，在 $\eta=0$ 时退化为 $SU(2)$ 线性态 SU2Bn0 的拟设。经过傅里叶变换，我们发现自旋子谱由下式决定：
\begin{equation*}
H=-2\chi\sin(k_x)\Gamma_0+2\eta\cos(k_x)\Gamma_2-2\chi\sin(k_y)\Gamma_1+2\eta\cos(k_y)\Gamma_3
\end{equation*}
其中 $k_x\in(-\pi/2,\pi/2)$，$k_y\in(-\pi,\pi)$，而
\begin{equation*}
\Gamma_0=\tau^{0}\otimes\tau^{3},\qquad
\Gamma_2=\tau^{1}\otimes\tau^{3},\qquad
\Gamma_1=\tau^{0}\otimes\tau^{1},\qquad
\Gamma_3=\tau^{2}\otimes\tau^{1}.
\end{equation*}
自旋子谱可以精确计算，它的四支具有 $\pm E_1(\pmb{k})$ 与 $\pm E_2(\pmb{k})$ 的形式。自旋子能量在两个孤立点 $\pmb{k}=(0,0)$ 与 $(0,\pi)$ 处为零。在 $\pmb{k}=0$ 附近，低能谱由下式给出（见图 9.10(b)）
\begin{equation*}
E=\pm\eta^{-1}\sqrt{(\chi^{2}+\eta^{2})^{2}(k_x^{2}-k_y^{2})^{2}+4\chi^{4}k_x^{2}k_y^{2}}
\end{equation*}
有趣（而且奇怪）的是，当 $\pmb{k}\to 0$ 时能量并不线性地趋于零，而是像 $k^{2}$ 那样趋于零！为了强调这种 $E\propto k^{2}$ 的二次方色散，我们将把这样的态称为 $Z_2$ 二次方自旋液体（$Z_2$-quadratic spin liquid）。

\subsection*{问题 9.6.1}

拟设 (9.6.3) 的 PSG。
\begin{enumerate}
\item 试由标记 Z2A001n 求出 Z2A001n PSG 的规范变换 $G_x$、$G_y$、$G_{P_x}$ 等。
\item 证明拟设 (9.6.3) 在 Z2A001n PSG 下不变。
\item 试求把 Z2A001n PSG 变换到 Z2A003n PSG 的规范变换。
\end{enumerate}

\subsection*{问题 9.6.2}

验证：对于由式 (9.6.10) 描述的自旋液体，只要 $\chi$ 与 $\eta$ 都不为零，绕回路 $\pmb{i}\to\pmb{i}+\pmb{x}\to\pmb{i}+\pmb{x}+\pmb{y}\to\pmb{i}+\pmb{y}\to\pmb{i}$ 与 $\pmb{i}\to\pmb{i}+\pmb{y}\to\pmb{i}-\pmb{x}+\pmb{y}\to\pmb{i}-\pmb{x}\to\pmb{i}$ 的 $SU(2)$ 通量算符（见式 (9.2.20)）就互不对易。因此，该自旋液体确实具有 $Z_2$ 规范结构。

\subsection*{问题 9.6.3}

(a) 证明：当 $\eta=0$ 时，式 (9.6.10) 中的拟设退化为 $SU(2)$ 线性态 SU2Bn0 的拟设；当 $\chi=0$ 时，退化为 $SU(2)$ 无能隙态 SU2An0 的拟设。（译注：原书此句将 $\eta$、$\chi$ 与两个态的对应关系印反，且把 SU2Bn0 误标为 “$SU(2)$-gapless”，现按 9.6.2 节正文及物理内容订正。）

(b) 试求自旋子谱 $\pm E_1$ 与 $\pm E_2$。

\section{量子序的物理测量}

\begin{keypoints}
\item 用 PSG 描述的量子序可以通过激发谱来测量。
\end{keypoints}

在用 PSG 从数学上刻画了量子序之后，我们想问：如何在实验中测量量子序？有能隙态中的量子序与拓扑序相联系，我们可以用基态简并、边缘态、准粒子统计等手段来测量拓扑序（Wen, 1990, 1995; Senthil and Fisher, 2001）。具有无能隙激发的态中的量子序则需要用不同的方式来测量。在本节中，我们想证明：一般而言，量子序可以通过无能隙激发的动力学性质来测量。然而，并非所有的动力学性质都是普适的。因此，在用无能隙激发来刻画并测量量子序之前，我们需要先甄别出无能隙激发的普适性质。量子序的 PSG 刻画使我们能够得到这些普适性质——我们只需找出具有同一 PSG 的所有拟设的无能隙激发所共有的性质。

为了演示上述想法，我们想研究双自旋子激发的谱。我们注意到，自旋子只能成对地产生，因此单自旋子谱不是物理的。我们还注意到，双自旋子谱中包含自旋 1 激发，而自旋 1 激发可以用中子散射实验来测量。

在给定动量处，双自旋子谱分布在一段或几段能量范围内。设 $E_{2s}(\pmb{k})$ 为动量 $\pmb{k}$ 处双自旋子谱的低边缘。在平均场理论中，双自旋子谱可以由单自旋子色散按如下方式构造：
\begin{equation*}
E_{2\text{-spinon}}(\pmb{k})=E_{1\text{-spinon}}(\pmb{q})+E_{1\text{-spinon}}(\pmb{k}-\pmb{q})
\end{equation*}
在图 9.11--9.14 中，我们给出了几个简单自旋液体的平均场 $E_{2s}$。如果平均场态对规范涨落是稳定的，我们预计平均场 $E_{2s}$ 应当与真实的 $E_{2s}$ 在定性上一致。

在我们的例子中，有八个稳定的 $Z_2$ 线性自旋液体（见图 9.11 与 9.12）。以这八个 $Z_2$ 自旋液体为例，我们可以演示自旋 1 激发的普适性质如何区分不同的量子序。

%FIG{9.11}: {$Z_2$ 线性自旋液体的 $E_{2s}(\pmb{k})$ 作为 $(k_x/2\pi,\,k_y/2\pi)$ 的函数的等高线图：(a) 式 (9.6.1) 中的 Z2A0013 态；(b) 式 (9.6.2) 中的 Z2Azz13 态；(c) 式 (9.6.3) 中的 Z2A001n 态；(d) 式 (9.6.4) 中的 Z2Azz1n 态。}

(a) 通过考察自旋 1 激发的谱，可以把 Z2A 自旋液体与 Z2B 自旋液体区分开。对于 Z2B 自旋液体，自旋 1 谱以布里渊区的四分之一为周期（即谱在 $k_x\to k_x+\pi$ 与 $k_y\to k_y+\pi$ 下不变）。与之相反，Z2A 自旋液体中的自旋 1 谱不具有这种周期性。

(b) 自旋 1 谱的周期性还可以把 Z2A0013 与 Z2Azz13 自旋液体同 Z2A001n 与 Z2Azz1n 自旋液体区分开。对于 Z2A001n 与 Z2Azz1n 自旋液体，自旋 1 谱以布里渊区的二分之一为周期（即谱在 $(k_x,k_y)\to(k_x+\pi,k_y+\pi)$ 下不变）。此外，Z2A001n 与 Z2Azz1n 自旋液体恰好在 $(\pi,\pi)$、$(\pi,0)$ 和 $(0,\pi)$ 处有无能隙自旋 1 激发。

(c) Z2A0013 与 Z2Azz13 自旋液体可以通过考察 $(\pi,0)$ 和 $(0,\pi)$ 附近的无能隙自旋 1 激发来区分。Z2Azz13 自旋液体中的无能隙自旋 1 激发位于区边界上。

%FIG{9.12}: {$Z_2$ 线性自旋液体的 $E_{2s}(\pmb{k})$ 作为 $(k_x/2\pi,\,k_y/2\pi)$ 的函数的等高线图：(a) 式 (9.6.5) 中的 Z2B0013 态；(b) 式 (9.6.6) 中的 Z2Bzz13 态；(c) 式 (9.6.7) 中的 Z2B001n 态；(d) 式 (9.6.8) 中的 Z2Bzz1n 态。谱以布里渊区的四分之一为周期。}

我们看到，自旋液体中的量子序可以通过探测自旋 1 激发的中子散射实验来测量。

接下来，让我们讨论 $U(1)$ 线性态 U1Cn01n（交错通量态）。人们曾提出用 U1Cn01n 态描述欠掺杂高 $T_c$ 超导体中的赝能隙金属态（Wen and Lee, 1996; Rantner and Wen, 2001）。U1Cn01n 态自然地解释了欠掺杂金属态中的赝能隙。作为一种代数自旋液体，U1Cn01n 态还解释了赝能隙态中的类 Luttinger 电子谱函数（Rantner and Wen, 2001; Franz and Tesanovic, 2001）以及 $(\pi,\pi)$ 自旋涨落的增强（Kim and Lee, 1999; Rantner and Wen, 2002）。

%FIG{9.13}: {$E_{2s}(\pmb{k})$ 作为 $(k_x/2\pi,\,k_y/2\pi)$ 的函数的等高线图：(a) 式 (9.2.39) 中的 $Z_2$ 无能隙态 Z2Ax2(12)n；(b) 式 (9.6.10) 中的 $Z_2$ 二次方态 Z2Bx2(12)n。}

%FIG{9.14}: {两个 $U(1)$ 线性自旋液体的 $E_{2s}(\pmb{k})$ 作为 $(k_x/2\pi,\,k_y/2\pi)$ 的函数的等高线图：(a) 式 (9.2.34) 中的 U1Cn01n 态（交错通量相）；(b) 无能隙相中式 (9.8.9) 的 U1Cn00x 态。谱以布里渊区的二分之一为周期。}

从图 9.14 我们看到，U1Cn01n 态中自旋 1 激发的无能隙点总是位于 $\pmb{k}=(\pi,\pi)$、$(0,0)$、$(\pi,0)$ 和 $(0,\pi)$。在所有这四个 $\pmb{k}$ 点处，自旋 1 连续谱边缘的等能线都具有两个彼此交叠的椭圆的形状。而且，等能线不垂直于区边界。所有这些都是 U1Cn01n 态的普适性质。在中子散射实验中测量这些性质，将使我们能够判断赝能隙金属态是否由 U1Cn01n（交错通量）态描述。

由于 $(2+1)$ 维 $U(1)$ 规范理论的瞬子效应，U1Cn01n 态是不稳定的。因此，U1Cn01n 态不得不变为其他一些态，例如 9.6 节中讨论的八个 $Z_2$ 自旋液体，或这里未讨论的其他态。从图 9.11(a) 我们看到，如果观察到 $(\pi,\pi)$ 处的节点劈裂为 $(\pi\pm\delta,\pi\pm\delta)$ 处的四个节点，以及 $(\pi,0)$ 和 $(0,\pi)$ 处的节点劈裂为 $(\pi\pm\delta,0)$ 和 $(0,\pi\pm\delta)$ 处的两个节点，就可以用中子散射探测从 U1Cn01n 态到 $Z_2$ 线性态 Z2A0013 的相变。从图 9.11(b) 我们看到，对于从 U1Cn01n 态到 $Z_2$ 线性态 Z2Azz13 的相变，$(\pi,\pi)$ 处的节点仍劈裂为 $(\pi\pm\delta,\pi\pm\delta)$ 处的四个节点，但 $(\pi,0)$ 和 $(0,\pi)$ 处的节点以另一种方式劈裂为 $(\pi,\pm\delta)$ 和 $(\pm\delta,\pi)$ 处的两个节点。我们还可以研究从 U1Cn01n 态到其余六个 $Z_2$ 自旋液体的相变。我们发现，自旋 1 激发谱以某些具有特征的方式改变。因此，通过测量自旋 1 激发谱及其演化，我们不仅能够探测不改变任何对称性的量子相变，还能够辨别正在发生的是哪一个相变。

\section{大 $N$ 极限下 $J_1$--$J_2$ 模型的相图}

\subsection{大 $N$ 极限}

\begin{keypoints}
\item 自旋-1/2 模型可以推广为 $SP(2N)$ 模型。
\item 对于 $SP(2N)$ 模型，由 $SU(2)$ 投影构造得到的平均场理论涨落微弱，是一个好的近似。
\end{keypoints}

到目前为止，我们一直集中讨论如何刻画、分类和测量不同的自旋液体及其量子序，还没有讨论如何寻找物理自旋哈密顿量，来实现我们构造的数百个不同自旋液体中的一些。本节就要处理这个问题。

在平均场层次上，设计一个实现具有给定量子序的自旋液体的自旋哈密顿量并不很难；对一个给定的自旋哈密顿量求出其平均场基态也不很难。真正的问题是我们是否应当相信平均场结果。我们已经看到，如果所得到的平均场基态是不稳定的（即如果平均场涨落导致低能下发散的相互作用），那么平均场结果不可信，平均场态不对应任何真实的物理自旋态。我们也曾论证过，如果平均场基态是稳定的（即平均场涨落导致低能下消失的相互作用），那么平均场结果可信，平均场态确实对应一个真实的物理自旋态。

这里我们想指出，上面关于稳定平均场态的论述过于乐观。一个“稳定平均场态”在低能下没有发散的涨落，所以它不一定不稳定；但另一方面，它也不一定稳定。这是因为短距离涨落如果足够强，同样会导致相变和失稳。因此，要使一个平均场结果可靠，平均场态必须稳定（或边缘），\emph{而且}短距离涨落必须微弱。由于我们的自旋模型中没有任何小参数，即便是稳定的平均场态，其短距离涨落也不微弱。正因如此，即使是稳定态的平均场结果能否应用于我们的自旋-1/2 模型，也并不清楚。

下面我们把自旋模型推广到大 $N$ 模型。我们将证明，大 $N$ 模型的平均场态具有微弱的短距离涨落，因此大 $N$ 模型的稳定与边缘平均场态对应于真实的物理态，稳定态或边缘态的平均场结果可以应用于大 $N$ 模型。

让我们从 $SU(2)$ 平均场理论的如下路径积分表示出发：
\begin{equation*}
\begin{aligned}
Z &= \int \mathcal{D}\psi\,\mathcal{D}[a_0^l(\pmb{i})]\,\mathcal{D}U_{ij}\ \mathrm{e}^{\mathrm{i}\int\mathrm{d}t\,(L-\sum_i a_0^l(\pmb{i},t)\psi_i^{\dagger}\tau^{l}\psi_i)} \\
L &= \sum_i \psi_i^{\dagger}\mathrm{i}\partial_t\psi_i-\sum_{\langle ij\rangle}\frac{1}{4}J_{ij}\left[\frac{1}{2}\mathrm{Tr}\,U_{ij}U_{ij}^{\dagger}-(\psi_i^{\dagger}U_{ij}\psi_j+h.c)\right].
\end{aligned}
\tag{9.8.1}
\end{equation*}
它由 $SU(2)$ 平均场哈密顿量 (9.2.9) 得到，其中 $U_{ij}$ 具有式 (9.2.13) 给出的形式。\footnote{我们把系数从 $3/8$ 改成了 $1/4$，这样，在式 (9.8.1) 中把 $U_{ij}$ 积掉就会得到自旋哈密顿量 (9.1.1)。}把费米子积掉之后，我们得到 $U_{ij}$ 与 $a_0^l$ 的有效拉格朗日量：
\begin{equation*}
Z=\int \mathcal{D}[a_0^l(\pmb{i})]\,\mathcal{D}U_{ij}\ \mathrm{e}^{\mathrm{i}\int\mathrm{d}t\,L_{0,eff}(U_{ij},a_0^l)}
\end{equation*}
问题在于，在上述路径积分中，$U_{ij}$ 与 $a_0^l$ 的涨落并不微弱。

为了减小涨落，我们引入 $N$ 份费米子副本 $\psi_i^a$，把式 (9.8.1) 推广为如下路径积分（Ran and Wen, 2003）：
\begin{equation*}
\begin{aligned}
Z &= \int \mathcal{D}\psi\,\mathcal{D}[a_0^l(\pmb{i})]\,\mathcal{D}U_{ij}\ \mathrm{e}^{\mathrm{i}\int\mathrm{d}t\,(L-\sum_i a_0^l(\pmb{i},t)\psi_i^{a\dagger}\tau^{l}\psi_i^{a})} \\
L &= \sum_i \psi_i^{a\dagger}\mathrm{i}\partial_t\psi_i^{a}-\sum_{\langle ij\rangle}\frac{1}{4}J_{ij}\left[\frac{N}{2}\mathrm{Tr}\,U_{ij}U_{ij}^{\dagger}-(\psi_i^{a\dagger}U_{ij}\psi_j^{a}+h.c)\right]
\end{aligned}
\tag{9.8.2}
\end{equation*}
其中 $a=1,2,\dots,N$。把费米子积掉之后，我们得到
\begin{equation*}
Z=\int \mathcal{D}[a_0^l(\pmb{i})]\,\mathcal{D}U_{ij}\ \mathrm{e}^{\mathrm{i}\int\mathrm{d}t\,N L_{0,eff}(U_{ij},a_0^l)}
\end{equation*}
我们看到，在大 $N$ 极限下，$U_{ij}$ 与 $a_0^l$ 的涨落是微弱的，平均场近似是一个好的近似。同样清楚的是，大 $N$ 平均场理论是一个 $SU(2)$ 规范理论，$U_{ij}$ 与 $a_0^l$ 的涨落对应于 $SU(2)$ 规范涨落。

在 9.2.1 节中，我们从物理自旋哈密顿量出发构造了平均场哈密顿量。这里我们面临的是相反的问题：已知大 $N$ 平均场哈密顿量
\begin{equation*}
H_{\text{mean}}=\sum_{\langle ij\rangle}\frac{1}{4}J_{ij}\left[\frac{N}{2}\mathrm{Tr}\,U_{ij}U_{ij}^{\dagger}-(\psi_i^{a\dagger}U_{ij}\psi_j^{a}+h.c)\right]+\sum_i a_0^l\psi_i^{a\dagger}\tau^{l}\psi_i^{a}
\tag{9.8.3}
\end{equation*}
我们想求出与之对应的物理的大 $N$ 自旋模型。

构造物理模型最重要的一步是找到物理希尔伯特空间。物理希尔伯特空间是费米子希尔伯特空间的一个子空间，由 $SU(2)$ 规范不变态、即满足约束
\begin{equation*}
\psi_i^{a\dagger}\tau^{l}\psi_i^{a}|phy\rangle,\qquad l=1,2,3
\end{equation*}
的态构成，且该约束在每个格点 $\pmb{i}$ 上都要成立。

得到每个格点上的物理希尔伯特空间之后，我们还需要找到在物理希尔伯特空间内部起作用的物理算符。这些物理算符是 $SU(2)$ 规范不变算符（即与 $\psi_i^{a\dagger}\tau^{l}\psi_i^{a}$ 对易的算符）。让我们把每个格点上 $\psi$ 的所有 $SU(2)$ 规范不变双线性型写下来：
\begin{equation*}
\begin{aligned}
S^{ab+}&\equiv\frac{1}{2}\psi_\alpha^{a\dagger}\tilde\psi_\alpha^{b},\qquad
S^{ab-}\equiv\frac{1}{2}\tilde\psi_\alpha^{a\dagger}\psi_\alpha^{b},\\
S^{ab3}&\equiv\frac{1}{2}\left(\psi_\alpha^{a\dagger}\psi_\alpha^{b}-\delta^{ab}\right)
=\frac{1}{2}\left(\delta^{ab}-\tilde\psi_\alpha^{b\dagger}\tilde\psi_\alpha^{a}\right)
\end{aligned}
\end{equation*}
其中 $\tilde\psi^{a}\equiv\mathrm{i}\sigma_2\psi^{a*}$，且格点指标已经隐去。这些 $S$ 算符是自旋-1/2 模型的自旋算符的推广。

当 $N=1$ 时，大 $N$ 模型就成为自旋-1/2 模型，$S$ 算符生成 $SU(2)$ 自旋旋转群。当 $N>1$ 时，$S$ 算符生成的是什么群？首先，让我们数一数不同 $S$ 算符的个数。对于 $S^{ab+}$ 或 $S^{ab-}$，标记关于 $a$ 与 $b$ 是对称的，因此每种类型各有 $N(N+1)/2$ 个不同的算符。对于 $S^{ab3}$，两个标记并不对称，因此就简单地是 $N^{2}$ 个。总共有 $N(N+1)+N^{2}=2N^{2}+N$ 个不同的 $S$ 算符。

可以考察 $S$ 算符之间的如下对易关系：
\begin{equation*}
\begin{gathered}
\left[S^{ab3},S^{cd3}\right]=\frac{1}{2}\left(\delta^{bc}S^{ad3}-\delta^{ad}S^{cb3}\right)\\[2pt]
\left[S^{ab3},S^{cd+}\right]=\frac{1}{2}\left(\delta^{bc}S^{ad+}+\delta^{bd}S^{ac+}\right)\\[2pt]
\left[S^{ab3},S^{cd-}\right]=-\frac{1}{2}\left(\delta^{ad}S^{bc-}+\delta^{ac}S^{bd-}\right)\\[2pt]
\left[S^{ab+},S^{cd-}\right]=\frac{1}{2}\left(\delta^{ac}S^{bd3}+\delta^{ad}S^{bc3}+\delta^{bc}S^{ad3}+\delta^{bd}S^{ac3}\right)\\[2pt]
\left[S^{ab-},S^{cd-}\right]=0,\qquad \left[S^{ab+},S^{cd+}\right]=0
\end{gathered}
\end{equation*}
这些正是 $SP(2N)$ 代数的关系式。所以，$S$ 算符就是生成 $SP(2N)$ 群的那 $2N^{2}+N$ 个生成元。当 $N=1$ 时，$SP(2)$ 同构于 $SU(2)$ 自旋旋转群。

在式 (9.8.2) 中把 $U_{ij}$ 与 $a_0^l$ 积掉，我们得到大 $N$ 模型的如下物理哈密顿量：
\begin{equation*}
H=\sum_{\langle ij\rangle}\frac{J_{ij}}{N}S_i^{ab}\cdot S_j^{ba}
\end{equation*}
其中
\begin{equation*}
\begin{aligned}
S_i^{ab1}&=S_i^{ba1}=\frac{1}{2}\left(S_i^{ab+}+S_i^{ab-}\right)\\
S_i^{ab2}&=S_i^{ba2}=\frac{1}{2\mathrm{i}}\left(S_i^{ab+}-S_i^{ab-}\right),
\end{aligned}
\end{equation*}
而
\begin{equation*}
S_i^{ab}=\left(S_i^{ab1},S_i^{ab2},S_i^{ab3}\right)
\end{equation*}
我们发现 $H$ 与 $SP(2N)$ 生成元 $\sum_i S_i^{ab}$ 对易，因此我们的大 $N$ 模型具有 $SP(2N)$ 对称性。这里我们想指出，$S_i^{ab}$ 的三个分量实际上并非处于同等地位：前两个关于 $a$ 与 $b$ 标记是对称的，第三个则不然。

我们知道，当 $N=1$ 时，物理希尔伯特空间在每个格点上有两个态，这两个态构成 $SP(2)=SU(2)$ 的一个不可约表示。对更高的 $N$，每格点物理希尔伯特空间的维数为
\begin{equation*}
\begin{aligned}
N &=\ 2\ \ \ 3\ \ \ 4\ \ \ 5\ \ \ 6\ \ \dots\\
\text{Dimension} &=\ 5\ \ \ 14\ \ \ 43\ \ \ 142\ \ \ 429\ \ \dots
\end{aligned}
\tag{9.8.4}
\end{equation*}
结果发现，这些物理希尔伯特空间是 $SP(2N)$ 对称群的不可约表示。我们可以用不可约表示在某个特定 Cartan 基下的最高权态来标记它。$SP(2N)$ 的 Cartan 基可以取为每个 $\alpha$ 的 $z$ 分量自旋，即
\begin{equation*}
S^{\alpha\alpha 3},\qquad \alpha=1,2,\dots,N.
\tag{9.8.5}
\end{equation*}
于是，物理希尔伯特空间中的最高权态就是不含 $\psi$ 费米子的态。

\subsection{$SP(2N)$ 模型的相图}

\begin{keypoints}
\item 平均场相图就是大 $N$ 极限下 $SP(2N)$ 模型的相图。
\item 平均场理论或 $SP(2N)$ 模型中不同的相可以用 PSG 来标记。
\end{keypoints}

从式 (9.8.3) 我们看到，$SP(2N)$ 模型的 $SU(2)$ 平均场理论与自旋-1/2 模型的 $SU(2)$ 平均场理论具有相同的结构。对两个模型而言，平均场态都由拟设 $(U_{ij},a_0^l(\pmb{i}))$ 描述。特别地，$SP(2N)$ 模型的平均场能量恰好是 $SP(2)$ 模型（即自旋-1/2 模型）平均场能量的 $N$ 倍。因此，平均场相图不依赖于 $N$。正如前面讨论的自旋-1/2 模型那样，$SP(2N)$ 模型的不同平均场态用 PSG 来刻画和分类。

这里我们考虑正方格子上一个具体的 $SP(2N)$ 模型。该模型具有最近邻耦合 $J_1$ 和次近邻耦合 $J_2$。我们取 $J_1+J_2=1$。

%FIG{9.15}: {$J_1$--$J_2$ 自旋系统中各个相的平均场能量：“A” 标记 $\pi$ 通量态（$SU(2)$ 线性态 SU2Bn0）；“B” 标记式 (9.8.6) 中的 $SU(2)\times SU(2)$ 无能隙态；“C” 标记式 (9.8.7) 中的 $SU(2)\times SU(2)$ 线性态；“D” 标记手征自旋态（一个 $SU(2)$ 有能隙态）；“E” 标记破坏 $90^{\circ}$ 旋转对称性的 $U(1)$ 线性态 (9.8.8)；“F” 标记式 (9.8.9) 中的 $U(1)$ 有能隙态 U1Cn00x；“G” 标记式 (9.6.2) 中的 $Z_2$ 线性态 Z2Azz13；“H” 标记式 (9.6.1) 中的 $Z_2$ 线性态 Z2A0013；“I” 标记均匀 RVB 态（$SU(2)$ 无能隙态 SU2An0）。}

在大 $N$ 极限下，$SU(2)$ 平均场理论 (9.8.3) 是一个好的近似，因此我们将用平均场理论来计算 $SP(2N)$ 模型的相图。描述平均场基态的平均场拟设通过把平均场能量极小化而求得，结果示于图 9.15。图 9.15 中的每条曲线代表一个局域极小的平均场能量，能量最低的曲线对应于平均场基态。我们发现，每条曲线上不同的拟设都具有同一个 PSG。这在预料之中，因为当我们改变 $J_2$ 时，曲线所描述的平均场能量以解析的方式变化，所以沿同一条曲线从一个点移动到另一个点并不发生量子相变。同一条曲线上的拟设属于同一个相，由同一个 PSG 描述。

然而，如果两条曲线相交，交点就代表一个量子相变，因为基态能量在交点处不是解析的。由于不同的曲线具有不同的 PSG，我们看到量子相变以 PSG 的改变为特征。下面我们将讨论图 9.15 中各平均场相的不同量子序（或 PSG）。

在图 9.15 中，相 A 是式 (9.2.33) 给出的 $\pi$ 通量态（SU2Bn0 $SU(2)$ 线性态）。相 B 是在对角链上带有两个相互独立的均匀 RVB 态的态。它在低能下具有 $SU(2)\times SU(2)$ 规范涨落，将被称为 $SU(2)\times SU(2)$ 无能隙态。其拟设由下式给出：
\begin{equation*}
u_{\pmb{i},\pmb{i}+\pmb{x}+\pmb{y}}=\chi\tau^{3},\qquad
u_{\pmb{i},\pmb{i}+\pmb{x}-\pmb{y}}=\chi\tau^{3},\qquad
a_0^{l}=0.
\tag{9.8.6}
\end{equation*}
相 C 是在对角链上带有两个相互独立的 $\pi$ 通量态的态。它在低能下具有 $SU(2)\times SU(2)$ 规范涨落，将被称为 $SU(2)\times SU(2)$ 线性态。其拟设由下式给出：
\begin{equation*}
u_{\pmb{i},\pmb{i}+\pmb{x}+\pmb{y}}=\chi(\tau^{3}+\tau^{1}),\qquad
u_{\pmb{i},\pmb{i}+\pmb{x}-\pmb{y}}=\chi(\tau^{3}-\tau^{1}),\qquad
a_0^{l}=0
\tag{9.8.7}
\end{equation*}
相 D 是手征自旋态 (9.1.22)。相 E 由拟设
\begin{equation*}
\begin{aligned}
u_{\pmb{i},\pmb{i}+\pmb{x}+\pmb{y}}&=\chi_1\tau^{1}+\chi_2\tau^{2}, &\qquad
u_{\pmb{i},\pmb{i}+\pmb{x}-\pmb{y}}&=\chi_1\tau^{1}-\chi_2\tau^{2}\\
u_{\pmb{i},\pmb{i}+\pmb{y}}&=\eta\tau^{3}, &\qquad
a_0^{l}&=0
\end{aligned}
\tag{9.8.8}
\end{equation*}
描述，它破坏 $90^{\circ}$ 旋转对称性，是一个 $U(1)$ 线性态。相 F 由如下 U1Cn00x 拟设描述：
\begin{equation*}
\begin{aligned}
u_{\pmb{i},\pmb{i}+\pmb{x}}&=\eta\tau^{1}, &\qquad
u_{\pmb{i},\pmb{i}+\pmb{y}}&=\eta\tau^{1}\\
u_{\pmb{i},\pmb{i}+\pmb{x}+\pmb{y}}&=\chi\tau^{3}, &\qquad
u_{\pmb{i},\pmb{i}+\pmb{x}-\pmb{y}}&=\chi\tau^{3}\\
a_0^{3}&=\lambda, &\qquad
a_0^{1,2}&=0.
\end{aligned}
\tag{9.8.9}
\end{equation*}
U1Cn00x 态可以是 $U(1)$ 线性态（若 $a_0^{3}$ 小），也可以是 $U(1)$ 有能隙态（若 $a_0^{3}$ 大）。相 F 的态结果是一个 $U(1)$ 有能隙态。相 G 由式 (9.6.2) 中的 Z2Azz13 拟设描述，是一个 $Z_2$ 线性态。相 H 由式 (9.6.1) 中的 Z2A0013 拟设描述，也是一个 $Z_2$ 线性态。相 I 是均匀 RVB 态（式 (9.2.32) 中的 $SU(2)$ 无能隙态 SU2An0）。

从图 9.15 我们观察到下列各对相之间（在平均场层次上）的连续相变：(A,D)、(A,G)、(B,G)、(C,E) 与 (B,H)。其中三个连续相变 (B,G)、(B,H) 与 (A,G) 不改变任何对称性。我们还注意到，在从相 A 到相 G 的连续相变中，相 A 的 $SU(2)$ 规范结构破缺到 $Z_2$；在 (B,G) 与 (B,H) 这两个相变中，相 B 的 $SU(2)\times SU(2)$ 规范结构破缺到 $Z_2$。

\section{量子序与平均场自旋液体的稳定性}

\begin{keypoints}
\item 许多无能隙平均场自旋液体对量子涨落可以是稳定的。即使存在长程规范相互作用，它们仍可以是稳定的。
\end{keypoints}

我们已经强调了平均场态对平均场涨落的稳定性的重要性。只有稳定的平均场态和边缘的平均场态才有可能描述物理自旋液体。我们在 9.1.5 节和 9.2.2 节中讨论过获得稳定平均场态的途径。这里我们将从量子序及其 PSG 刻画的角度再次讨论平均场态的稳定性。
